% !TeX root = ../../tesis.tex
%%%%%%%%%%%%%%%%%%%%%%%%%%%%%%%%%%%%%%%%%%%%%%%%%%%%%%%%%%%%%%%%%%%
% SECCIÓN 5.5.2 — CENTRALIDAD DE INTERMEDIACIÓN B(v)
%%%%%%%%%%%%%%%%%%%%%%%%%%%%%%%%%%%%%%%%%%%%%%%%%%%%%%%%%%%%%%%%%%%
\subsection{Centralidad de Intermediación (\texorpdfstring{$B(v)$}{B(v)})}
\label{subsec:fase3-centralidad}

No todos los nodos de la red tienen el mismo peso en la conectividad
global. Algunos funcionan como puentes inevitables: la mayoría de las
rutas mínimas entre cualquier par de la red pasan por ellos, y su
ausencia fragmentaría la estructura o alargaría los tiempos de forma
generalizada. La centralidad de intermediación $B(v)$ mide exactamente
esa condición: la proporción de rutas mínimas del grafo que atraviesan
cada nodo.

Dado un grafo $\mathcal{G}$ con nodos $\mathcal{V}$, la centralidad de
intermediación de un nodo $v$ se define como:

\begin{equation*}
    B(v) = \sum_{s \neq v \neq t} \frac{\sigma_{st}(v)}{\sigma_{st}}
    \tag{\ref{eq:betweenness}}
\end{equation*}

\noindent donde:
\begin{itemize}
    \item $\sigma_{st}$: cantidad de rutas mínimas entre los nodos $s$ y $t$.
    \item $\sigma_{st}(v)$: subconjunto de esas rutas que pasan por $v$.
\end{itemize}

El modelo calcula $B(v)$ mediante el algoritmo de Brandes, que reduce la
complejidad computacional a $\mathcal{O}(VE)$ sobre grafos ponderados; la
corrida sobre el componente gigante (10,561~nodos) tomó 593.4~segundos.

Dicho de otro modo: hay estaciones que sostienen la red desde adentro
sin que los usuarios lo noten. No son necesariamente las más
transitadas; son las que aparecen en la ruta cuando no hay otra
opción. Quitarlas no molesta a unos cuantos pasajeros —desorganiza
trayectos que atraviesan cuadrantes enteros de la metrópoli.

El Cuadro~\ref{tab:bv-top10} presenta los diez nodos con mayor $B(v)$
calculados sobre el componente gigante. El resultado más llamativo:
los diez son estaciones del Metro, y Tacubaya encabeza la lista con
un valor de 0.2179.

\begin{table}[H]
    \centering
    \caption[Top 10 nodos por centralidad de intermediación $B(v)$]{%
    Diez nodos con mayor centralidad de intermediación ($B(v)$) en la
    red de transporte de la \zmvm{}. Corrida: 2026-09-02.}
    \label{tab:bv-top10}
    \small
    \begin{tabular}{|r|l|l|l|r|}
        \hline
        \textbf{\#} & \textbf{Estación} & \textbf{Sistema} &
        \textbf{Alcaldía} & \textbf{$B(v)$} \\
        \hline
         1 & Tacubaya        & Metro & Miguel Hidalgo & 0.2179 \\
         2 & Mixcoac         & Metro & Benito Juárez  & 0.1725 \\
         3 & Hidalgo         & Metro & Cuauhtémoc     & 0.1595 \\
         4 & Balderas        & Metro & Cuauhtémoc     & 0.1433 \\
         5 & Lázaro Cárdenas & Metro & Cuauhtémoc     & 0.1376 \\
         6 & Chilpancingo    & Metro & Cuauhtémoc     & 0.1366 \\
         7 & Viaducto        & Metro & Benito Juárez  & 0.1346 \\
         8 & Villa de Cortés & Metro & Benito Juárez  & 0.1341 \\
         9 & Nativitas       & Metro & Benito Juárez  & 0.1299 \\
        10 & Zapata          & Metro & Benito Juárez  & 0.1283 \\
        \hline
    \end{tabular}
    \fuentefigura{Elaboración propia con \texttt{VFTModel}
    \citep{cabrera2026VFTModel}. Corrida: 2026-09-02.}
\end{table}

El Metro monopoliza los diez primeros puestos pese a concentrar
menos del 2\,\% de las aristas de servicio real: su ventaja no es
volumétrica sino posicional ---sus estaciones ocupan los puntos donde
más rutas se cruzan.

La distribución geográfica del top~10 es igualmente estrecha: siete
de los diez nodos quedan en Benito Juárez o Cuauhtémoc, lo que revela
que la intermediación de la red está concentrada en un corredor
reducido del centro-poniente de la cuenca.

La Figura~\ref{fig:grafica-bv-nodos} extiende ese ranking a los
principales nodos de la red.

\begin{figure}[H]
    \centering
    \includegraphics[width=0.92\textwidth]{%
        Figures/Mapas/Tableau_Resources/RedActual/VFTModel/UniGrafica/%
        Grafica_VFTModel_2026_Centralidad_Intermediacion_Nodos}
    \caption[Centralidad de intermediación $B(v)$ --- top nodos de la red]{%
    Ranking de los nodos con mayor centralidad de intermediación
    ($B(v)$) en la red de transporte de la \zmvm{}.
    Corrida: 2026-09-02.}
    \label{fig:grafica-bv-nodos}
    \fuentefigura{Elaboración propia con \textbf{Tableau};
    datos: \texttt{VFTModel} \citep{cabrera2026VFTModel}.}
\end{figure}

La centralidad topológica no siempre coincide con la afluencia real.
Para verificarlo, la Figura~\ref{fig:grafica-top10-estaciones} muestra
las diez estaciones con mayor volumen de usuarios según los registros
de Apímetro ---un contraste directo con el ranking de $B(v)$ del
Cuadro~\ref{tab:bv-top10}.

\begin{figure}[H]
    \centering
    \includegraphics[width=0.92\textwidth]{%
        Figures/Mapas/Tableau_Resources/RedActual/Apimetro/UniGrafica/%
        Grafica_Apimetro_2026_Top_10_estaciones}
    \caption[Top 10 estaciones por afluencia de usuarios --- \zmvm{}]{%
    Diez estaciones con mayor volumen de usuarios registrado por
    \texttt{Apímetro} en la red de transporte de la \zmvm{}.
    Corrida: 2026-04-18.}
    \label{fig:grafica-top10-estaciones}
    \fuentefigura{Elaboración propia con \textbf{Tableau};
    datos: \texttt{Apímetro} \citep{cabrera2025apimetro}.}
\end{figure}

La Figura~\ref{fig:grafica-bv-sistema} agrega esos valores por sistema,
mostrando cuánto peso acumulado concentra cada modo en la intermediación
global de la red.

\begin{figure}[H]
    \centering
    \includegraphics[width=0.92\textwidth]{%
        Figures/Mapas/Tableau_Resources/RedActual/VFTModel/UniGrafica/%
        Grafica_VFTModel_2026_Centralidad_Intermediacion_sistema}
    \caption[Centralidad de intermediación $B(v)$ acumulada por sistema]{%
    Suma acumulada de $B(v)$ agrupada por sistema de transporte.
    El Metro concentra la mayor proporción de intermediación global
    de la red. Corrida: 2026-09-02.}
    \label{fig:grafica-bv-sistema}
    \fuentefigura{Elaboración propia con \textbf{Tableau};
    datos: \texttt{VFTModel} \citep{cabrera2026VFTModel}.}
\end{figure}

Para traducirlo al lenguaje urbano: imaginemos que cada par de colonias
de la \zmvm{} traza su ruta más rápida hacia el destino. Las estaciones
que aparecen con mayor frecuencia en ese ejercicio colectivo no son
necesariamente las más concurridas en términos de afluencia, sino las
más difíciles de sustituir: su ausencia obligaría a miles de trayectos
a buscar desvíos más largos o, en el peor caso, a no encontrarlos.
$B(v)$ mide exactamente esa indispensabilidad estructural.

La Figura~\ref{fig:mapa-betweenness} traslada ese ranking al territorio,
mostrando cómo la intermediación topológica se distribuye espacialmente
sobre la cuenca metropolitana.

\begin{figure}[H]
    \centering
    \includegraphics[width=0.92\textwidth]{%
        Figures/Mapas/Qgis_Resources/RedActual/Mapa_Betweenness}
    \caption[Mapa de centralidad de intermediación $B(v)$ --- \zmvm{}]{%
    Distribución espacial de la centralidad de intermediación ($B(v)$)
    en la red de transporte de la \zmvm{}. El corredor de alta
    intermediación traza una banda en el centro-poniente de la cuenca.
    Corrida: 2026-09-02.}
    \label{fig:mapa-betweenness}
    \fuentefigura{Elaboración propia con \textbf{QGIS};
    datos: \texttt{VFTModel} \citep{cabrera2026VFTModel}.}
\end{figure}

La concentración que revelan los datos tabulares adquiere forma geográfica
en el mapa: el corredor de alta intermediación no cubre la metrópoli de
manera uniforme, sino que traza una banda relativamente estrecha sobre el
centro-poniente, dejando amplias zonas periféricas con valores bajos o
nulos. Esas zonas ---las mismas que registraron menor cobertura y mayor
Factor de Desviación en Fase~1--- son también las que dependen de los
nodos del corredor central para cualquier viaje de largo aliento.

La Figura~\ref{fig:dashboard-bv} consolida la visión integral del indicador
sobre el conjunto de la red.

\begin{figure}[H]
    \centering
    \includegraphics[width=0.92\textwidth]{%
        Figures/Mapas/Tableau_Resources/RedActual/VFTModel/Dashboards/%
        Dashboard_VFTModel_2026_Centralidad_intermediacion}
    \caption[Dashboard integral de centralidad de intermediación $B(v)$ --- \zmvm{}]{%
    Vista integral del indicador de centralidad de intermediación ($B(v)$)
    sobre la red de transporte de la \zmvm{}.
    Corrida: 2026-09-02.}
    \label{fig:dashboard-bv}
    \fuentefigura{Elaboración propia con \textbf{Tableau};
    datos: \texttt{VFTModel} \citep{cabrera2026VFTModel}.}
\end{figure}

Los tres indicadores de Fase~3 ---tiempo de viaje, intermediación y
vulnerabilidad preliminar--- apuntan en la misma dirección: la red
concentra su capacidad de articulación en un corredor reducido del
centro-poniente, mientras la periferia sur y oriente queda
estructuralmente subordinada a él. La sección siguiente examina qué
ocurre cuando alguno de esos nodos críticos desaparece de la red.
